\chapter{本体与大语言模型}
\chapterart{ch08}{本体作为大模型输出的「守门人」}

\begin{introbox}
\textbf{【导读】}
大语言模型能流畅地谈论任何话题——包括它并不真正「知道」的事。
当它被问到一台不在训练数据中的设备参数时，往往会给出一个语气笃定、数值具体、
却完全凭空生成的答案，这就是\textbf{幻觉}（hallucination）。
本章给出幻觉的工程解法：不指望靠训练或提示词消除它，
而是用本体在系统层面\textbf{拦截}它。
我们将沿四条主线展开：三层闸门架构（生成前注入约束、生成中强制查询通道、
生成后事实校验）、Text2SPARQL 的安全转换工程、GraphRAG 的结构化检索，
以及给智能体装上本体护栏的完整协议。最后讨论方向相反的问题——
大模型反过来如何加速本体工程本身，以及整套系统该如何评测。
本章是全书与 AI 工程衔接最紧密的一章，也回应了第 1 章的开篇之问：
为什么 AI 时代的企业必须建立受来源约束的语义权威层；该层不替代现实事实提供者和有权决策者。
\end{introbox}

\chapterbinding{semantica.chapter\_packages.vol1.ch08}{partial}{blocked}
{幻觉控制、GraphRAG、Text2SPARQL 与 legacy Agent 案例已经迁入；
场景 \texttt{OE-V1-CH08-SCN-GUARDRAIL-001} 保留参考 verdict。}
{其操作合同状态为 \texttt{blocked\_unsupported\_operation\_contract}；
模型、prompt、ontology snapshot、工具动作与授权尚未端到端绑定，历史 Python 案例不得冒充执行收据。}

\section{幻觉的机理：为什么大模型会一本正经地胡说}

先看一个会让任何工艺工程师心头一紧的场景。某车间上线了一个对话式工艺助手，
新来的操作员问它：「Lathe\_003 加工钛合金时主轴转速应该设多少？」
助手流畅作答：「建议设置为 1800 转/分，并将进给速度控制在 0.15 毫米/转。」
数字具体、语气专业、格式规范——唯一的问题是，这台车床根本没有加工钛合金的能力认证，
而那两个「参数」是模型从训练语料中无数台别的机床的描述里「平均」出来的。
如果操作员照做，轻则废一批料，重则损伤主轴。

\subsection{统计生成的本质}

要理解幻觉为什么必然出现，只需要理解大语言模型在做什么：
给定前文，预测下一个词的概率分布，取样输出，再把输出拼回前文，如此循环。
模型的全部「知识」都压缩在权重里，而权重的唯一来源是训练语料的统计规律。
这个机制有两个直接推论。

第一，\textbf{模型没有「事实承诺」}。它生成「15.5 kW」和生成「22 kW」
走的是同一条概率采样路径，没有任何内部机制区分「我查证过」与「我觉得像」。
流畅与正确在统计生成里是两个独立维度——训练优化的是前者，业务需要的是后者。

第二，\textbf{分布外必然失真}。你工厂里那台 Lathe\_003 的真实参数从未出现在
训练语料中，模型对它的一切「了解」都是从「车床」这个词的全局统计中外推的。
外推在语言风格上几乎总是成功（这正是它显得可信的原因），
在事实层面则只能靠运气。

\subsection{幻觉的四种典型形态}

工程场景中观察到的幻觉可以归为四类，识别形态有助于设计针对性的拦截手段：

\begin{center}
\begin{tabularx}{\linewidth}{@{}lXX@{}}
\toprule
\textbf{类型} & \textbf{表现（制造场景例）} & \textbf{危险点} \\
\midrule
凭空捏造 & 为不存在的设备编出型号、参数、维护记录 & 数值具体，最难凭直觉识破 \\
张冠李戴 & 把铣床的参数安到车床上；把 A 区设备说成在 B 区 & 单看每个事实都「真」，组合错误 \\
过度泛化 & 「所有数控车床都能加工钛合金」 & 类级断言听起来像专业知识 \\
时效错位 & 引用已废止的旧工艺规程、已停产的型号 & 内容曾经为真，难以用常识排除 \\
\bottomrule
\end{tabularx}
\end{center}

值得注意的是「张冠李戴」：它的每个碎片都来自真实语料，
拼接错误发生在生成过程中，因此任何「让模型多读真实资料」的方案都治不了它——
拼接这个动作本身就是统计性的。

\subsection{为什么提示词工程治标不治本}

一个常见的误解是：在系统提示里写上「请不要编造，不知道就说不知道」就能解决幻觉。
实践中这类指令确实能降低幻觉频率——模型会变得更保守、更多使用模糊限定词——
但它改变的是模型的\textbf{表达倾向}，不是模型的\textbf{知识状态}。
模型并不知道自己不知道：「Lathe\_003 的功率」在它的内部表征里
与「车床的典型功率」纠缠在一起，它无法区分哪部分是针对这台具体设备的可靠记忆，
哪部分是类别统计的外推。让一个无法自知边界的系统自觉守住边界，逻辑上就不成立。

因此工程上的结论是：提示词约束值得做（它是第一层闸门的一部分），
但\textbf{可信度必须由模型之外的系统层提供}。这正是本体的位置。

\subsection{代价分级：什么场景必须设防}

并非所有幻觉都值得动用全套防线。按输出的去向分级，可以合理分配工程投入：
闲聊与头脑风暴场景，幻觉代价近似为零，无需设防；
内部草稿场景（写报告初稿、整理会议纪要），有人工复核兜底，轻度设防即可；
对外答复场景（客服、技术支持），错误直接损害信任与合规，必须全量校验；
\textbf{驱动动作}场景（Agent 修改系统状态、下达指令），错误会物化为物理世界的损失，
除全量校验外还需要审计与回滚能力——这是 8.5 节的主题。

\section{三层闸门：把可信度从模型内搬到系统层}

既然模型内部无法自证可靠，解决思路就清晰了：
在模型的上游、中游、下游各设一道闸门，让\textbf{事实只能来自本体与知识图谱}，
让违反领域公理的内容根本无法流出系统。先看全景：

\begin{codebox}
\codeline{\textcolor{cmtgray}{\#\ ==============================}}
\codeline{\textcolor{cmtgray}{\#\ 2.\ 三层控制架构}}
\codeline{\textcolor{cmtgray}{\#\ ==============================}}
\codeline{\textcolor{cmtgray}{\#}}
\codeline{\textcolor{cmtgray}{\#\ \ \ 用户问题}}
\codeline{\textcolor{cmtgray}{\#\ \ \ \ \ │}}
\codeline{\textcolor{cmtgray}{\#\ \ \ [第一层\ 生成前]\ 约束注入（Prompt阶段）}}
\codeline{\textcolor{cmtgray}{\#\ \ \ \ \ │\ \ \ ·\ 把本体词汇表/类定义注入上下文："只能使用以下设备类型……"}}
\codeline{\textcolor{cmtgray}{\#\ \ \ \ \ │\ \ \ ·\ 把取值范围注入："功率单位kW，范围(0,10000]"}}
\codeline{\textcolor{cmtgray}{\#\ \ \ \ \ │\ \ \ ·\ 声明边界："知识库中没有的信息回答'无记录'"}}
\codeline{\textcolor{cmtgray}{\#\ \ \ \ \ │}}
\codeline{\textcolor{cmtgray}{\#\ \ \ [第二层\ 生成中]\ 工具调用（受控通道）}}
\codeline{\textcolor{cmtgray}{\#\ \ \ \ \ │\ \ \ ·\ LLM不直接回答事实，而是生成查询（SPARQL/函数调用）}}
\codeline{\textcolor{cmtgray}{\#\ \ \ \ \ │\ \ \ ·\ 事实只能来自知识图谱返回值（见\ text2sparql.txt）}}
\codeline{\textcolor{cmtgray}{\#\ \ \ \ \ │}}
\codeline{\textcolor{cmtgray}{\#\ \ \ [第三层\ 生成后]\ 事实校验（出口闸门）}}
\codeline{\textcolor{cmtgray}{\#\ \ \ \ \ │\ \ \ ·\ 实体校验：回答中的实体能否链接到本体个体？}}
\codeline{\textcolor{cmtgray}{\#\ \ \ \ \ │\ \ \ ·\ 断言校验：回答中的三元组是否存在于图谱/可由推理导出？}}
\codeline{\textcolor{cmtgray}{\#\ \ \ \ \ │\ \ \ ·\ 范围校验：数值是否通过SHACL形状检查？}}
\codeline{\textcolor{cmtgray}{\#\ \ \ \ \ │\ \ \ ·\ 一致性校验：是否违反不相交/函数性公理？}}
\codeline{\textcolor{cmtgray}{\#\ \ \ \ \ ▼}}
\codeline{\textcolor{cmtgray}{\#\ \ \ 通过\ \(\rightarrow\)\ 输出（附溯源）；不通过\ \(\rightarrow\)\ 拒答或降级为"无记录"}}
\codeblank{}
\end{codebox}

\assetsource{ch08}{hallucination-control}

三层闸门各司其职，也各有边界。下面逐层展开。

\subsection{生成前：约束注入}

第一层在提示词阶段工作，向模型注入三类信息。
\textbf{词汇表注入}：把本体的类与属性清单（或与当前问题相关的子集）写进上下文，
并声明「只能使用以下词汇描述设备与能力」。
这利用了大模型出色的指令跟随能力——它无法保证不幻觉，
但能显著收窄输出的词汇空间，让后续校验更容易命中。
\textbf{取值范围注入}：把关键数值属性的物理区间写进上下文
（「功率单位 kW，合法区间 (0, 10000]」），
给模型的数值生成提供锚点。
\textbf{边界声明}：明确告知「知识库中没有的信息，回答『无记录』，不要推测」，
并给出「无记录」的标准话术模板。

第一层的价值是「降低后续闸门的压力」，它的边界也很清楚：
上下文长度有限，无法注入完整本体；模型对指令的遵守是概率性的。
所以它永远不能单独使用。

\subsection{生成中：受控通道}

第二层是整个架构的支柱：\textbf{让 LLM 不直接回答事实问题，
而是生成对知识图谱的查询}。用户问「3 号车床现在什么状态」，
模型的产出不是「空闲」这两个字，而是一条 SPARQL：
查询发给图谱端点，返回的结果才是事实，模型再把结果组织成自然语言。

这个设计的深意在于改变了信息流的拓扑：原本是
「训练语料 \(\rightarrow\) 模型记忆 \(\rightarrow\) 回答」，
现在变成「用户问题 \(\rightarrow\) 模型翻译 \(\rightarrow\)
图谱检索 \(\rightarrow\) 回答」。模型从「知识的来源」降级为「语言的翻译器」，
而翻译错误（查询写错）会在执行时显式失败，远比「编一个流畅的答案」容易发现。
第二层的完整工程展开见 8.3 节（Text2SPARQL）。

\subsection{生成后：出口校验}

第三层假设前两层都可能漏过，对最终输出做四项独立校验。
\textbf{实体校验}：回答中提到的每个命名实体（设备、材料、订单号）
必须能链接到本体中的个体，链接不上的实体视为捏造。
\textbf{断言校验}：从回答中抽取「主语—谓语—宾语」形式的事实断言，
逐条用 ASK 查询核验它是否存在于图谱中、或能否由推理导出。
\textbf{范围校验}：数值断言走一遍 SHACL 形状检查（第 7 章），
功率为负、日期在未来这类「物理不可能」被直接拦截。
\textbf{一致性校验}：把候选断言临时加入知识库，跑一次推理器一致性检查——
违反不相交、函数性等公理的断言会让知识库不一致，该输出即被拒绝。

四项校验中任何一项不通过，系统的选择不是「修一修再发」，
而是\textbf{拒答或降级}：要么明确告知「该问题无法可靠回答」，
要么把可校验的部分输出、把存疑的部分替换为「无记录」。
宁可少说，不可说错——这是对外系统的底线策略。

\subsection{完整对照：同一个问题的两种命运}

把三层闸门串起来，看包内教学资产中的完整对照：

\begin{codebox}
\codeline{\textcolor{cmtgray}{\#\ 问："Lathe\_003\ 能加工钛合金吗？它功率多大？"}}
\codeblank{}
\codeline{\textcolor{cmtgray}{\#\ 无控制（直接生成）：}}
\codeline{\textcolor{cmtgray}{\#\ \ \ "Lathe\_003功率22kW，可以加工钛合金。"\ \ ←\ 两个事实都是编造}}
\codeblank{}
\codeline{\textcolor{cmtgray}{\#\ 三层控制下：}}
\codeline{\textcolor{cmtgray}{\#\ \ \ 第一层注入：本体中\ Equipment/Material\ 词汇\ +\ canProcess\ 定义}}
\codeline{\textcolor{cmtgray}{\#\ \ \ 第二层查询：}}
\codeline{\textcolor{cmtgray}{\#\ \ \ \ \ ASK\ \{\ mfg:Lathe\_003\ mfg:canProcess\ mfg:TitaniumAlloy\ \}\ \ \(\rightarrow\)\ false}}
\codeline{\textcolor{cmtgray}{\#\ \ \ \ \ SELECT\ ?p\ \{\ mfg:Lathe\_003\ mfg:hasPower\ ?p\ \}\ \ \ \ \ \ \ \ \ \ \ \ \ \(\rightarrow\)\ 15.5}}
\codeline{\textcolor{cmtgray}{\#\ \ \ 第三层校验后输出：}}
\codeline{\textcolor{cmtgray}{\#\ \ \ \ \ "功率15.5kW（来源：设备台账）。知识库中没有该设备可加工}}
\codeline{\textcolor{cmtgray}{\#\ \ \ \ \ \ 钛合金的能力记录（注意：'无记录'不等于'不能'，}}
\codeline{\textcolor{cmtgray}{\#\ \ \ \ \ \ 如需确认请联系工艺部门）。"}}
\codeblank{}
\end{codebox}
\color{black}

\assetsource{ch08}{hallucination-control}

逐步拆解路径 B 发生了什么：第一层注入了设备与材料的词汇表，
使模型在翻译问题时使用规范术语；第二层把两个事实问题
（能力、功率）各翻译成一条查询，ASK 返回 false、SELECT 返回 15.5；
第三层校验输出时确认 15.5 落在 SHACL 区间内、
「无加工钛合金记录」的表述不与任何公理冲突。
最终回答的每个成分都有出处，而且——注意最后一句——
「无记录」与「不能」被严格区分，这是第 2 章开放世界假设在产品话术上的直接落地。

\subsection{这不是专家系统的还魂}

读到这里，熟悉 AI 历史的读者可能心生疑问：
规则、公理、推理器——这不就是上世纪专家系统那一套吗？
形似而神异，差别在三处。
\textbf{语言入口换了}：专家系统的死穴是知识获取瓶颈与僵硬的交互
（用户必须按系统的语言提问）；现在 LLM 承担了全部语言理解，
用户随便问，系统内部才切换到形式化轨道。
\textbf{知识层标准化了}：当年的规则库是各家私有格式，
今天的本体是 W3C 标准栈（第 4 章），跨系统、跨厂商、可交换。
\textbf{分工哲学变了}：专家系统试图用规则覆盖一切智能，
本章的架构只让公理负责「守底线」——开放性的理解、综合、表达全部交给模型。
一句话：不是用规则替代神经网络，而是\textbf{用一小撮硬规则
给一个强大但不可靠的神经网络划出安全区}。
这个组合的两端都是上一代技术做不到的。

\subsection{一次请求的完整时序走查}

把三层闸门串成时间线，跟踪「3 号车床能加工钛合金吗？功率多大？」
这条请求的完整旅程，能把抽象架构落到毫秒级的直觉上。
（1）请求进入，会话层取出对话实体栈——空栈，无指代需要解析。
（2）实体链接：「3 号车床」命中别名表 \(\rightarrow\) \texttt{mfg:Lathe\_003}，
「钛合金」命中 \(\rightarrow\) \texttt{mfg:TitaniumAlloy}；毫秒级。
（3）第一层注入：按链接结果选取本体子集（设备与材料相关的类、属性、范围）
拼装系统提示；毫秒级。
（4）意图识别：能力查询 + 参数查询，两个子问题；
模板库命中「能力查询」模板，参数查询走生成路线。
（5）生成路线产出 SPARQL，进入四道校验：语法通过、
SELECT 在白名单、\texttt{mfg:hasPower} 在 Schema 中、自动附加 LIMIT；
生成是全链路延迟大头，秒级。
（6）两条查询打到只读端点执行：ASK 返回 false，SELECT 返回 15.5；
图谱查询数十毫秒。
（7）结果交回模型组织回答；秒级。
（8）第三层出口校验：回答中的实体可链接、断言「功率 15.5」与图谱一致、
SHACL 区间通过、无公理冲突；亚秒级。
（9）话术检查：无记录部分套用标准话术。
（10）PROV 记录落库：查询原文、结果、来源图、时间戳。
（11）回答送出。（12）对话实体栈压入 \texttt{Lathe\_003}，等待下一轮。
全程两次模型调用、两次图谱查询，端到端三到五秒——
这就是「可信」在工程上的真实成本结构：大头是模型推理，
而所有保证可信的校验环节加起来只占零头。

\subsection{性价比最高的一条公理}

三层闸门中，公理校验的拦截能力取决于本体里写了多少「硬约束」。
经验表明，\textbf{不相交公理}（disjointness）是性价比最高的一类：

\begin{codebox}
\codeline{\textcolor{cmtgray}{\#\ ==============================}}
\codeline{\textcolor{cmtgray}{\#\ 4.\ 类型混淆的公理拦截}}
\codeline{\textcolor{cmtgray}{\#\ ==============================}}
\codeline{\textcolor{cmtgray}{\#\ LLM常见错误：把概念归入错误类型（嫁接幻觉）}}
\codeline{\textcolor{cmtgray}{\#\ 防线：互不相交公理\ +\ 出口校验}}
\codeblank{}
\codeline{\textcolor{cmtgray}{\#\ 本体公理：}}
\codeline{\textcolor{cmtgray}{\#\ \ \ AllDisjointClasses(CNCLathe,\ CNCMillingMachine,\ Robot,\ AGV)}}
\codeline{\textcolor{cmtgray}{\#\ \ \ Equipment\ \(\sqcap\)\ Material\ \(\sqsubseteq\)\ \(\bot\)}}
\codeblank{}
\codeline{\textcolor{cmtgray}{\#\ 拦截示例：}}
\codeline{\textcolor{cmtgray}{\#\ \ \ LLM输出断言\ (Lathe\_003,\ rdf:type,\ Robot)}}
\codeline{\textcolor{cmtgray}{\#\ \ \ \(\rightarrow\)\ 与已有断言\ (Lathe\_003,\ rdf:type,\ CNCLathe)\ 在不相交公理下矛盾}}
\codeline{\textcolor{cmtgray}{\#\ \ \ \(\rightarrow\)\ 推理器报不一致\ \(\rightarrow\)\ 该输出被拒绝}}
\codeblank{}
\codeline{\textcolor{cmtgray}{\#\ 工程提示：不相交公理是性价比最高的防幻觉公理——}}
\codeline{\textcolor{cmtgray}{\#\ 建模时对每组同层兄弟类显式声明\ AllDisjointClasses}}
\codeblank{}
\end{codebox}

\assetsource{ch08}{hallucination-control}

大模型的「张冠李戴」型幻觉，本质上是把不同类别的属性混拼。
而建模时只要对每组同层兄弟类显式声明 \texttt{AllDisjointClasses}
（车床、铣床、机器人、AGV 两两不相交），
任何「把 Lathe\_003 当机器人调度」的断言都会在一致性检查中立刻现形。
这条建议落实成本极低——在 Protégé 中选中兄弟类点一下 Disjoint 即可——
强烈建议作为建模规范固定下来（呼应第 3 章的建模纪律）。

\subsection{留下证据链：PROV 审计}

闸门拦截解决「错的不出去」，审计解决「对的可追溯」：

\begin{codebox}
\codeline{\textcolor{cmtgray}{\#\ ==============================}}
\codeline{\textcolor{cmtgray}{\#\ 5.\ 溯源与审计（PROV思想）}}
\codeline{\textcolor{cmtgray}{\#\ ==============================}}
\codeline{\textcolor{cmtgray}{\#\ 每个对外输出附带证据链：}}
\codeline{\textcolor{cmtgray}{\#\ \ \ 结论：\ \ \ \ 功率15.5kW}}
\codeline{\textcolor{cmtgray}{\#\ \ \ 证据：\ \ \ \ (mfg:Lathe\_003,\ mfg:hasPower,\ 15.5)}}
\codeline{\textcolor{cmtgray}{\#\ \ \ 来源图：\ \ graph\ <http://example.org/graph/ledger>（台账，可信度高）}}
\codeline{\textcolor{cmtgray}{\#\ \ \ 查询：\ \ \ \ 所执行的SPARQL原文}}
\codeline{\textcolor{cmtgray}{\#\ \ \ 时间戳：\ \ 2026-06-11T10:30:00+08:00}}
\codeline{\textcolor{cmtgray}{\#\ 价值：事后可审计"AI为什么这么说"；}}
\codeline{\textcolor{cmtgray}{\#\ 高合规行业（航空DO-178C、汽车ISO\ 26262、医药GxP）可直接复用该审计链}}
\codeblank{}
\end{codebox}

\assetsource{ch08}{hallucination-control}

每个对外输出都携带证据四元组（结论、支撑三元组、来源图、执行的查询与时间戳），
事后任何人都能重放「AI 为什么这么说」。
在高合规行业（航空 DO-178C、汽车 ISO 26262、医药 GxP），
这条审计链不是锦上添花，而是系统能否获准上线的前提条件。
溯源的数据模型直接复用第 7 章的具名图与来源元数据——
基础设施做对了，审计是「顺便」得到的。

\subsection{拒答与可用性：一条绕不开的权衡曲线}

把闸门拧得越紧，幻觉漏过率越低，但「其实答得上来却被拒答」的比例也会升高。
这条权衡曲线没有免费午餐，工程上能做的是\textbf{把曲线整体推向右上}，
而不是在烂曲线上挑点。推曲线的三个抓手按收益排序：
扩充图谱覆盖（答案本来就在库里，拒答自然减少——这是第 7 章工作的直接回报）；
改进实体链接（大量「假拒答」其实是链接失败，别名表每补一行都在改善曲线）；
细化拒答话术（「无记录」与「问题不明确，您是指……吗」是两种不同的拒答，
后者把死路变成了对话）。
理解这条曲线还有一个管理学含义：给业务方汇报时，
拒答率上升未必是坏消息——如果同期幻觉漏过率下降更多，系统正在变得更诚实。

\subsection{三个常见误解的澄清}

实践中关于幻觉有三个流传很广的误解，逐一澄清能避免误投资源。

\textbf{误解一：「上了 RAG 就不会幻觉了」。}
检索增强确实让模型「有据可依」，但模型对检索结果的使用仍是统计性的：
它可能忽略检索内容、可能把两段检索结果错误拼接、
也可能在检索为空时悄悄退回参数记忆。RAG 改善的是「有没有材料」，
不解决「是否忠实于材料」——后者仍需要第三层闸门逐条核验。

\textbf{误解二：「把温度调到 0 就稳了」。}
温度为 0 只是让采样变成贪心选择，模型给出它认为概率最高的续写。
如果知识本身缺失，概率最高的续写依然是编造——只是每次编得一样而已。
确定性与正确性是两回事。

\textbf{误解三：「领域微调能解决」。}
微调能显著改善术语使用与回答风格，也能注入一批领域事实，
但事实会随业务漂移（设备新增、参数变更），而微调后的权重无法增量更新——
今天调好的模型，下个月就开始过期。对比之下，
知识图谱的更新是即时生效的（第 7 章的增量管道），
这正是「事实放图谱、语言放模型」分工的根本理由。

\subsection{三层闸门的失效模式与互补}

没有任何一层是完美的，架构的可靠性来自互补。逐层列出失效模式：

\begin{center}
\begin{tabularx}{\linewidth}{@{}lXX@{}}
\toprule
\textbf{闸门} & \textbf{典型失效} & \textbf{由谁兜底} \\
\midrule
生成前注入 & 模型忽略指令；上下文装不下完整词汇表 & 第二、三层 \\
受控通道 & 查询语义偏差（查了但查错）；实体链接错 & 第三层断言校验 \\
出口校验 & 断言抽取漏掉隐含陈述；校验集覆盖不全 & 审计抽样与红队（8.7） \\
\bottomrule
\end{tabularx}
\end{center}

注意第三层自己的失效要靠\textbf{流程}兜底而非更多代码：
生产抽样审计负责发现漏网之鱼，红队测试负责主动制造漏网场景。
安全工程的常识在这里同样适用——纵深防御的最后一环永远是运营。

\subsection{拒答与降级的话术工程}

第三层闸门做出「拦」的决定后，怎么对用户说同样需要设计。
两类标准话术值得做成模板沉淀下来。
\textbf{无记录话术}：「知识库中没有 Lathe\_003 可加工钛合金的能力记录
（截至 2026-06-10 的台账数据）。『无记录』不代表『不能』，
如需确认请联系工艺部门或提交能力认证申请。」——
三个成分缺一不可：明确说无记录、给出数据时点、区分无记录与否定并指出下一步。
\textbf{置信不足话术}：「您的问题可能涉及多台设备
（3 号车床 / 3 号铣床），请确认是哪一台。」——
把系统的不确定暴露为选择题而不是含糊其辞。
话术模板与查询模板一样进版本库管理；
它们是用户对系统「诚实程度」的直接感知面，值得产品与工程共同评审。

\subsection{三层闸门的最小落地路径}

全套架构看起来工程量不小，但它可以渐进落地，每一步都独立产生价值。
\textbf{第一周}：图谱端点切只读、上操作白名单——一行配置级别的工作，
先把「不可逆的错」锁死。
\textbf{第二到三周}：盘点高频问题做二十条查询模板 + 实体别名表，
模板法问答上线——此时已覆盖多数流量且零幻觉。
\textbf{第四到六周}：接入 LLM 生成路线与 Schema 校验，处理长尾问法；
同时上线断言 ASK 校验（第三层的最小版本）。
\textbf{之后}：SHACL 范围校验、一致性校验、PROV 审计、评测体系逐步补齐。
这条路径的设计原则是\textbf{先封死高危面、再扩展能力面}——
与安全工程「先关门、再开窗」的次序一致。

\subsection{部署形态}

三层闸门落到系统架构上有三种常见形态。
\textbf{网关式}：闸门做成独立服务，横在所有 LLM 应用与模型之间，
统一注入、统一校验——适合多应用共享一套知识治理的企业，
代价是引入一跳延迟与一个关键组件。
\textbf{库式}：闸门作为 SDK 嵌入各应用进程，延迟最低、
但版本碎片化风险高（各应用闸门版本不一致时，防线出现参差）。
\textbf{边车式}：闸门与应用同机部署、独立进程，
在隔离性与延迟之间折中，配合容器编排可以做到版本统一滚动升级。
选型的决定因素通常不是技术而是组织：谁对「AI 输出的可信度」负责，
闸门就应该部署在谁的控制边界内。

\section{Text2SPARQL：自然语言到查询的安全转换}

第二层闸门说「让模型写查询」，这一节回答怎么把这件事做到生产级。
表面上它只是一次翻译，实际上要同时解决三个问题：
翻译得对（语义保真）、翻译得安全（不被注入利用）、翻译失败时不乱来（回退有序）。

\subsection{三条技术路线}

\begin{codebox}
\codeline{\textcolor{cmtgray}{\#\ ==============================}}
\codeline{\textcolor{cmtgray}{\#\ 1.\ 三种技术路线}}
\codeline{\textcolor{cmtgray}{\#\ ==============================}}
\codeline{\textcolor{cmtgray}{\#\ 路线A\ 模板槽填充：}}
\codeline{\textcolor{cmtgray}{\#\ \ \ 预定义查询模板，LLM只负责识别意图与填槽}}
\codeline{\textcolor{cmtgray}{\#\ \ \ 高可控、零注入风险；覆盖面有限}}
\codeline{\textcolor{cmtgray}{\#\ 路线B\ LLM直接生成SPARQL\ +\ 校验层：}}
\codeline{\textcolor{cmtgray}{\#\ \ \ 灵活覆盖长尾问题；必须配校验层（见第3节）}}
\codeline{\textcolor{cmtgray}{\#\ 路线C\ 语义解析（训练专用模型）：}}
\codeline{\textcolor{cmtgray}{\#\ \ \ 精度高；需要标注数据，迭代成本高}}
\codeline{\textcolor{cmtgray}{\#\ 工程组合：高频问题走A，长尾走B，A/B共享同一校验层}}
\codeblank{}
\end{codebox}
\color{black}

\assetsource{ch08}{text2sparql}

三条路线不是竞争关系而是分工关系。
模板法把高频问题做成「填空题」，覆盖客服与车间问答中 60\%--80\% 的流量，
其确定性让它天然免疫注入与幻觉；
LLM 生成路线负责长尾的灵活问法，是表达力的来源，也是风险的来源，
必须配齐校验层；语义解析路线精度上限最高，但需要标注数据与持续训练，
适合查询模式稳定、预算充足的大型部署。
工程上的常见演化路径是：先用模板法上线拿到真实问句分布，
再用这些问句训练或调优生成路线，最后视规模决定是否引入专用解析模型。

\subsection{路线A精讲：模板槽填充}

\begin{codebox}
\codeline{\textcolor{cmtgray}{\#\ ==============================}}
\codeline{\textcolor{cmtgray}{\#\ 2.\ 路线A示例：模板槽填充}}
\codeline{\textcolor{cmtgray}{\#\ ==============================}}
\codeline{\textcolor{cmtgray}{\#\ 模板T01\ 设备状态查询：}}
\codeline{\textcolor{cmtgray}{\#\ \ \ 意图关键词：状态/在干嘛/空闲吗}}
\codeline{\textcolor{cmtgray}{\#\ \ \ 模板：}}
\codeline{\textcolor{cmtgray}{\#\ \ \ \ \ SELECT\ ?status\ WHERE\ \{\ \{EQUIPMENT\}\ mfg:hasStatus\ ?status\ \}}}
\codeline{\textcolor{cmtgray}{\#\ \ \ 槽位：\{EQUIPMENT\}\ ←\ 实体链接结果（别名表见第7章）}}
\codeblank{}
\codeline{\textcolor{cmtgray}{\#\ 模板T02\ 能力查询：}}
\codeline{\textcolor{cmtgray}{\#\ \ \ "哪些设备能加工\{MATERIAL\}？"}}
\codeline{\textcolor{cmtgray}{\#\ \ \ \ \ SELECT\ ?e\ WHERE\ \{}}
\codeline{\textcolor{cmtgray}{\#\ \ \ \ \ \ \ \ \ ?e\ rdf:type/rdfs:subClassOf*\ mfg:Equipment\ ;}}
\codeline{\textcolor{cmtgray}{\#\ \ \ \ \ \ \ \ \ \ \ \ mfg:canProcess\ \{MATERIAL\}\ .\ \}}}
\codeblank{}
\codeline{\textcolor{cmtgray}{\#\ 填槽过程：}}
\codeline{\textcolor{cmtgray}{\#\ \ \ "3号车床现在空闲吗"\ \(\rightarrow\)\ 意图T01\ +\ 实体mfg:Lathe\_003}}
\codeline{\textcolor{cmtgray}{\#\ \ \ \(\rightarrow\)\ SELECT\ ?status\ WHERE\ \{\ mfg:Lathe\_003\ mfg:hasStatus\ ?status\ \}}}
\codeline{\textcolor{cmtgray}{\#\ \ \ \(\rightarrow\)\ 返回\ mfg:Idle\ \(\rightarrow\)\ 生成回答"是，当前空闲"}}
\codeblank{}
\end{codebox}

\assetsource{ch08}{text2sparql}

模板法的两个环节都值得细看。
\textbf{意图识别}可以用关键词规则起步（示例中的「状态/在干嘛/空闲吗」），
量级上来后换成轻量分类模型；意图集合就是模板库的索引。
\textbf{实体链接}把口语化的「3 号车床」解析为 \texttt{mfg:Lathe\_003}——
注意这里直接复用了第 7 章实体对齐建好的别名表，
图谱工程的积累在这里第二次产生回报（第一次是融合数据时）。
槽位填充完成后，查询是从模板实例化出来的，
语法必然正确、操作必然只读，这就是模板法「天然安全」的含义。

\subsection{路线B的四道校验}

LLM 自由生成的查询则必须逐条过闸：

\begin{codebox}
\codeline{\textcolor{cmtgray}{\#\ ==============================}}
\codeline{\textcolor{cmtgray}{\#\ 3.\ 路线B的校验层（关键防线）}}
\codeline{\textcolor{cmtgray}{\#\ ==============================}}
\codeline{\textcolor{cmtgray}{\#\ LLM生成的SPARQL必须依次通过四道检查：}}
\codeblank{}
\codeline{\textcolor{cmtgray}{\#\ (1)\ 语法校验}}
\codeline{\textcolor{cmtgray}{\#\ \ \ \ \ 用SPARQL解析器（Jena\ ARQ\ /\ rdflib）解析，失败\(\rightarrow\)重试或回退模板}}
\codeblank{}
\codeline{\textcolor{cmtgray}{\#\ (2)\ 操作白名单（防注入/防破坏）}}
\codeline{\textcolor{cmtgray}{\#\ \ \ \ \ 只允许\ SELECT\ /\ ASK；}}
\codeline{\textcolor{cmtgray}{\#\ \ \ \ \ 拒绝\ INSERT\ /\ DELETE\ /\ DROP\ /\ LOAD\ /\ SERVICE（防外联）}}
\codeline{\textcolor{cmtgray}{\#\ \ \ \ \ 生产端点同时在服务端禁用update（双保险）}}
\codeblank{}
\codeline{\textcolor{cmtgray}{\#\ (3)\ Schema校验（防"词汇幻觉"）}}
\codeline{\textcolor{cmtgray}{\#\ \ \ \ \ 查询中出现的类/属性必须存在于本体：}}
\codeline{\textcolor{cmtgray}{\#\ \ \ \ \ \ \ LLM编造\ mfg:hasMaxSpeed（本体中只有\ mfg:hasSpindleSpeed）}}
\codeline{\textcolor{cmtgray}{\#\ \ \ \ \ \ \ \(\rightarrow\)\ Schema校验失败\ \(\rightarrow\)\ 把本体属性清单回喂LLM重新生成}}
\codeline{\textcolor{cmtgray}{\#\ \ \ \ \ 这是Text2SPARQL最常见错误，校验后准确率提升显著}}
\codeblank{}
\codeline{\textcolor{cmtgray}{\#\ (4)\ 资源限额}}
\codeline{\textcolor{cmtgray}{\#\ \ \ \ \ 强制\ LIMIT（如1000）、查询超时（如5s）、禁用笛卡尔积形态}}
\codeblank{}
\end{codebox}

\assetsource{ch08}{text2sparql}

该历史教学资产括号中列出的 Jena ARQ/rdflib 只是解析器生态举例，不是本书运行入口；
要升级为可执行主张，四道检查必须全部落在 ch08 的绑定 parser、snapshot、policy、
runner 与 oracle 内。当前 operation contract 缺失，因此本节仍是设计规格。
四道校验的次序经过设计，成本递增、收益分层。
\textbf{语法校验}最便宜，解析失败直接打回，连端点都不用碰。
\textbf{操作白名单}防的是注入：想象用户输入
「查一下 3 号车床状态，顺便清空所有维护记录」——
若模型照单全收生成了 DELETE 语句，白名单会在执行前拦下；
配合服务端只读端点（第 4 章提到的端点安全），形成双保险。
\textbf{Schema 校验}拦的是「词汇幻觉」：模型凭语感发明
\texttt{mfg:hasMaxSpeed} 这种本体中不存在的属性，是生成路线最高频的错误形态；
校验失败后把本体属性清单回喂给模型重试一次，命中率会显著回升。
实践经验是：四道校验中，Schema 校验对最终执行准确率的提升幅度最大——
它恰好补上了模型「不知道本体里有什么」这个结构性盲区。
\textbf{资源限额}（强制 LIMIT、超时、禁笛卡尔积形态）防的不是恶意而是失误，
保护端点不被一条写歪的查询拖垮。

\subsection{失败回退与评测}

任何一道校验失败，回退序列是：带着错误信息重试（最多两次）
\(\rightarrow\) 降级到模板匹配 \(\rightarrow\) 澄清式反问用户
（「您是想查状态还是维护记录？」）。

\begin{codebox}
\codeline{\textcolor{cmtgray}{\#\ ==============================}}
\codeline{\textcolor{cmtgray}{\#\ 5.\ 失败回退策略}}
\codeline{\textcolor{cmtgray}{\#\ ==============================}}
\codeline{\textcolor{cmtgray}{\#\ 生成失败/校验不过\ \(\rightarrow\)\ 重试（带错误信息）最多2次}}
\codeline{\textcolor{cmtgray}{\#\ \ \ \(\rightarrow\)\ 仍失败\ \(\rightarrow\)\ 回退模板匹配}}
\codeline{\textcolor{cmtgray}{\#\ \ \ \(\rightarrow\)\ 模板也无\ \(\rightarrow\)\ 澄清式反问用户（"您是想查状态还是维护记录？"）}}
\codeline{\textcolor{cmtgray}{\#\ 禁止：在任何失败分支让LLM"凭印象"直接回答事实}}
\codeblank{}
\end{codebox}

\assetsource{ch08}{text2sparql}

回退链的铁律写在最后一行：任何失败分支都不允许模型「凭印象」直接回答事实——
否则前面所有闸门等于虚设。
评测方面有一个一举两得的做法：第 3 章的 CQ 清单本来就是
「自然语言问法 \(\rightarrow\) 标准查询」的配对集，
直接把它扩充成 Text2SPARQL 的评测集（每个 CQ 配 3--5 种口语化变体），
同时度量语法正确率、Schema 命中率与执行准确率三层指标。

\subsection{生成路线的提示词构成}

路线B的输出质量，七成取决于喂给模型的上下文怎么组织。
一份生产级的 Text2SPARQL 系统提示包含五个要素，缺一不可。
\textbf{角色与边界}：「你是查询翻译器，只输出 SPARQL，不回答事实」——
一句话锁定任务形态，杜绝模型「顺手」直答。
\textbf{Schema 摘要}：类层次（缩进文本即可）加属性清单，
每个属性标注定义域与值域——这是防词汇幻觉的源头供给；
本体大时按实体链接结果动态选取相关子集。
\textbf{方言约定}：声明端点支持的 SPARQL 版本与扩展函数，
避免模型使用端点不认的语法。
\textbf{少样本示例}：三到五对「问题 \(\rightarrow\) 查询」，
覆盖单实体查询、过滤、聚合、属性路径四种形态——
示例的选择比数量重要，每个示例都应展示一种模型容易写错的模式。
\textbf{不确定出口}：「无法确定如何翻译时输出 NEED\_CLARIFY」——
给模型一条体面的退路，比让它硬写一条错查询便宜得多。

\subsection{实体链接的应用层细节}

实体链接在第 7 章解决的是「入库时两条记录是否同一实体」，
在这里解决的是「用户的一句话指向哪个个体」——同根同源，约束不同：
应用层链接必须在百毫秒内完成、必须容忍口语噪声、错了立刻被用户看见。
工程上分三步。\textbf{候选生成}：别名表精确命中优先；
未命中则在标签索引上做编辑距离与拼音模糊匹配，召回三到五个候选。
\textbf{消歧打分}：用上下文特征排序——对话实体栈中的近邻加分、
与问题中其他实体同车间/同工单的加分、设备类型与问题动词匹配
（「加工」倾向设备，「检验」倾向量具）加分。
\textbf{置信处理}：高分独占则直接采用；多个候选接近则反问用户
（呈现为选择题）；全部低分则按「实体幻觉」处理，触发无记录话术。
注意一条容易踩的坑：不要把消歧失败静默地「选第一个」——
宁可多一次反问，不要让用户拿着 3 号铣床的数据去操作 3 号车床。

\subsection{难点问法：聚合、时间与单位}

单实体查询翻译成功率最高，三类问法构成长尾中的硬骨头，值得专项打磨。
\textbf{聚合类}（「各工作单元的平均功率」）：模型常犯的错是
GROUP BY 变量与 SELECT 投影不匹配、或把 COUNT 用在了该用 AVG 的地方；
对策是在少样本示例里固定放一个聚合模板，并在 Schema 摘要中标注哪些属性可聚合。
\textbf{时间类}（「上个月的故障记录」）：「上个月」必须先被解析为
具体日期区间再进入查询——这个解析放在模型外做（系统时钟 + 日期库），
把确定的区间作为槽位喂给模型，比让模型自己算日期可靠得多
（呼应第 5 章：SWRL 没有 NOW，同样的问题同样的解法）。
\textbf{单位类}（「功率超过两万瓦的设备」）：图谱里存的是 kW，
问题里是瓦——单位归一化要在校验层显式处理，
否则一条语法完美、Schema 全中的查询会安静地返回错误答案。
这三类的共同教训是：\textbf{凡是有确定算法的子任务，
都从模型手里拿回来用代码做}，模型只负责真正需要语言理解的部分。

\subsection{多轮对话中的指代与上下文}

真实用户不会每句话都带全名：「3 号车床什么状态？」「那它今天的任务呢？」——
第二句的「它」需要解析回 \texttt{mfg:Lathe\_003}。
工程做法是维护一个轻量的\textbf{对话实体栈}：
每轮成功链接的实体压栈，后续轮次的代词与省略主语优先在栈顶解析；
解析结果作为槽位传给模板或写进生成路线的上下文。
栈深三到五层即可覆盖绝大多数对话，更深的引用（「我们最开始说的那台设备」）
让模型在 NEED\_CLARIFY 与候选列表之间选择，宁可反问不可猜。
这个组件同样复用第 7 章的别名表——到这里读者应该已经看清一个模式：
\textbf{图谱工程的每一分投入，都会在 LLM 应用层反复变现}。

\section{GraphRAG：让检索理解结构}

第二层闸门解决「事实从图谱来」，但当问题需要\textbf{综合}多个事实时，
还需要解决「把哪些事实给模型看」——这是检索增强生成（RAG）的领地。
本节讨论为什么朴素 RAG 在企业知识上力不从心，以及图谱如何补位。

\subsection{朴素RAG的两个结构性缺陷}

朴素 RAG 的流程是：文档切块、向量化入库，问题来了按语义相似度取回
top-k 个文本块塞进上下文。它在「答案集中在某一段文字里」的场景表现良好，
但企业知识有两个特征恰好打在它的软肋上。

\textbf{结构丢失}：切块把「设备—工序—产品」的关联切断了。
一个文本块写着「W001 工单包含车削与铣削工序」，
另一个块写着「车削工序需要数控车床」，
向量检索可能只召回其中一块——关系的另一半永远缺席。

\textbf{多跳失效}：问「P001 延期会影响哪些客户订单？」，
答案需要沿「产品 \(\rightarrow\) 工单 \(\rightarrow\) 订单 \(\rightarrow\) 客户」
走三跳。没有任何单一文本块同时包含整条链路，
相似度检索召回的每个块都只贴着问题的某个词，拼不出完整答案——
这不是参数没调好，是表示方式的上限。

\subsection{GraphRAG 的五步流程}

GraphRAG 的解法是换检索对象：不检索文本块，检索\textbf{图谱子图}。

\begin{codebox}
\codeline{\textcolor{cmtgray}{\#\ ==============================}}
\codeline{\textcolor{cmtgray}{\#\ 2.\ GraphRAG\ 流程}}
\codeline{\textcolor{cmtgray}{\#\ ==============================}}
\codeline{\textcolor{cmtgray}{\#\ \ \ 用户问题}}
\codeline{\textcolor{cmtgray}{\#\ \ \ \ \ │}}
\codeline{\textcolor{cmtgray}{\#\ \ \ (1)\ 实体识别与链接}}
\codeline{\textcolor{cmtgray}{\#\ \ \ \ \ │\ \ \ \ \ "P001延期影响哪些订单"\ \(\rightarrow\)\ 锚点实体\ mfg:P001}}
\codeline{\textcolor{cmtgray}{\#\ \ \ \ \ │}}
\codeline{\textcolor{cmtgray}{\#\ \ \ (2)\ 子图扩展（k-hop\ /\ 按关系类型走）}}
\codeline{\textcolor{cmtgray}{\#\ \ \ \ \ │\ \ \ \ \ 从P001出发沿\ producedBy/contains/orderedBy\ 扩展2\textasciitilde{}3跳}}
\codeline{\textcolor{cmtgray}{\#\ \ \ \ \ │\ \ \ \ \ 得到：P001\ ←\ W001(工单)\ ←\ O-2026-088(订单)\ ←\ 客户C17}}
\codeline{\textcolor{cmtgray}{\#\ \ \ \ \ │}}
\codeline{\textcolor{cmtgray}{\#\ \ \ (3)\ 子图裁剪与排序}}
\codeline{\textcolor{cmtgray}{\#\ \ \ \ \ │\ \ \ \ \ 按与问题的相关度裁剪（边类型白名单、路径打分），}}
\codeline{\textcolor{cmtgray}{\#\ \ \ \ \ │\ \ \ \ \ 控制在LLM上下文预算内}}
\codeline{\textcolor{cmtgray}{\#\ \ \ \ \ │}}
\codeline{\textcolor{cmtgray}{\#\ \ \ (4)\ 子图序列化\ \(\rightarrow\)\ 注入上下文}}
\codeline{\textcolor{cmtgray}{\#\ \ \ \ \ │}}
\codeline{\textcolor{cmtgray}{\#\ \ \ (5)\ LLM\ 基于子图作答（要求引用图中三元组作为依据）}}
\codeblank{}
\end{codebox}
\color{black}

\assetsource{ch08}{graphrag-pattern}

五步中有三个工程关键点。
\textbf{实体锚定}（第 1 步）复用 Text2SPARQL 的实体链接组件，
把问题中的「P001」「延期」「订单」映射为图谱锚点。
\textbf{子图扩展}（第 2 步）从锚点沿边扩展 k 跳，工程上 k 一般不超过 3：
跳数预算既是延迟控制也是噪声控制。
\textbf{裁剪排序}（第 3 步）按「边类型与问题的相关度」打分裁剪，
把子图压进上下文预算。

本体在这里再次扮演「轨道」角色：扩展时\textbf{只走本体里声明过的属性}
（关系白名单直接取自属性定义），节点按类层次过滤
（问订单就只保留 \texttt{Order} 及其子类的节点）。
没有本体约束的图谱扩展会带回大量噪声边；
有本体约束的扩展结果天然干净——这是「GraphRAG 因本体而可靠」的含义。

\subsection{裁剪打分的实做}

「按相关度裁剪」说起来轻巧，实做时打分函数通常由三个因子相乘或加权：
\textbf{边类型先验}——本体里每类属性预设一个与业务重要度对应的基础权重
（影响分析场景里 \texttt{partOf}、\texttt{orderedBy} 权重高，
\texttt{hasColor} 这类描述性属性权重低），这张权重表本身挂在本体上维护；
\textbf{问题相关度}——边或节点的标签与问题文本的语义相似度
（用同一套向量模型计算，复用已有组件）；
\textbf{结构惩罚}——高度数节点（连接成千上万边的「枢纽」，
比如「车间A」）的扩展按度数衰减，防止子图被枢纽节点炸开。
打分后按预算贪心保留最高分的边，直到节点数或 token 预算触顶。
经验上，权重表的初版交给领域专家拍十分钟即可，
后续靠评测集的失败案例微调——又一个「先拍脑袋、再用数据修」的务实循环。

\subsection{子图怎么喂给模型：序列化}

拿到子图后还有最后一个设计决策——以什么格式写进上下文：

\begin{codebox}
\codeline{\textcolor{cmtgray}{\#\ ==============================}}
\codeline{\textcolor{cmtgray}{\#\ 3.\ 子图序列化格式对比}}
\codeline{\textcolor{cmtgray}{\#\ ==============================}}
\codeline{\textcolor{cmtgray}{\#\ 格式A：三元组列表（最忠实，token较省）}}
\codeline{\textcolor{cmtgray}{\#\ \ \ (P001,\ producedBy,\ W001)}}
\codeline{\textcolor{cmtgray}{\#\ \ \ (W001,\ partOf,\ O-2026-088)}}
\codeline{\textcolor{cmtgray}{\#\ \ \ (O-2026-088,\ orderedBy,\ Customer\_C17)}}
\codeline{\textcolor{cmtgray}{\#\ \ \ (O-2026-088,\ dueDate,\ "2026-06-20")}}
\codeblank{}
\codeline{\textcolor{cmtgray}{\#\ 格式B：自然语言化（LLM理解最稳）}}
\codeline{\textcolor{cmtgray}{\#\ \ \ "产品P001由工单W001生产；W001属于订单O-2026-088；}}
\codeline{\textcolor{cmtgray}{\#\ \ \ \ 该订单的客户是C17，交期2026-06-20。"}}
\codeblank{}
\codeline{\textcolor{cmtgray}{\#\ 格式C：JSON（适合函数调用/结构化输出链路）}}
\codeline{\textcolor{cmtgray}{\#\ \ \ \{"entity":"P001","produced\_by":\{"workorder":"W001",}}
\codeline{\textcolor{cmtgray}{\#\ \ \ \ "order":\{"id":"O-2026-088","customer":"C17","due":"2026-06-20"\}\}\}}}
\codeblank{}
\codeline{\textcolor{cmtgray}{\#\ 实践：格式A+B混合——三元组保证可校验，旁附一句话说明降低误读}}
\codeblank{}
\end{codebox}
\color{black}

\assetsource{ch08}{graphrag-pattern}

三种格式的取舍是「忠实度、可读性、token 成本」的三角：
三元组列表最忠实且省 token，但模型偶尔会误读方向；
自然语言化最稳，token 开销最大；JSON 适合下游还要结构化处理的链路。
示例中给出的混合策略（三元组为主、旁附一句话说明）在实践中最常用——
三元组保证答案可以回溯校验（第三层闸门直接拿它去 ASK），
那句话则消除方向歧义。

\subsection{检索的时点一致性}

企业问题常带时间语境：「上周五排产时 3 号车床什么状态？」
如果检索总是返回\textbf{当前}图谱，答案对「现在」正确、对问题错误。
第 5 章与第 7 章铺设的具名图快照在这里闭环：
时间解析组件把「上周五排产时」转为时点，
检索层把查询路由到对应的快照图（或在支持时间戳边的存储上加时间过滤），
子图扩展与断言校验都在同一时点的切片上进行——
\textbf{检索、回答、校验三者的时间基准必须一致}，
否则会出现「答案是旧的、校验用新的、互相打架」的诡异故障。
对没有快照基础设施的轻量部署，至少要做到：
回答中显式标注数据时点（话术工程已包含），
让用户知道「这是截至何时的事实」。

\subsection{两个失败案例的剖析}

GraphRAG 的失败往往不是「检索不到」而是「检索回来一堆错的」，
两个典型案例值得引以为戒。
\textbf{案例一：没开白名单的扩展}。某团队上线初版时图省事，
从锚点沿\textbf{所有}边扩展两跳。问「P001 的质量问题」，
子图里混进了 P001 的包装规格、定价记录、推广素材——
因为图谱里这些边确实都连着 P001。模型面对一锅杂烩，
回答里混入了与质量毫不相干的「事实」，用户评价是「答非所问还特别自信」。
修复就是一张边类型白名单：质量场景只走
\texttt{hasDefect}、\texttt{producedBy}、\texttt{inspectedBy} 等六类边。
教训：\textbf{图谱越丰富，无差别扩展越危险}——丰富是资产，前提是检索有纪律。
\textbf{案例二：枢纽节点爆炸}。问「车间 A 上周的异常」，
「车间 A」这个节点连着数百台设备、数千条任务边，两跳扩展直接超出 token 预算，
裁剪后保留的边近乎随机。修复是结构惩罚（高度数节点降权）加
聚合改写——把「车间 A 的异常」改写为一条聚合查询（Text2SPARQL 路线）
而不是子图检索。教训呼应上一小节的协同判据：
\textbf{聚合形态的问题本来就不该走检索}。

\subsection{混合检索与社区摘要}

图检索与向量检索不是替代关系：

\begin{codebox}
\codeline{\textcolor{cmtgray}{\#\ ==============================}}
\codeline{\textcolor{cmtgray}{\#\ 4.\ 检索策略：向量\ vs\ 图\ vs\ 混合}}
\codeline{\textcolor{cmtgray}{\#\ ==============================}}
\codeline{\textcolor{cmtgray}{\#\ 向量检索：召回与问题语义相近的"描述性知识"（维修案例文本）}}
\codeline{\textcolor{cmtgray}{\#\ 图检索：\ \ 精确获取"结构性知识"（关联、路径、聚合）}}
\codeline{\textcolor{cmtgray}{\#\ 混合（推荐）：}}
\codeline{\textcolor{cmtgray}{\#\ \ \ 实体锚定\ \(\rightarrow\)\ 图扩展拿结构\ \(\rightarrow\)\ 对图中实体挂载的文档片段再做向量检索}}
\codeline{\textcolor{cmtgray}{\#\ \ \ （例：拿到BearingWear节点后，再向量检索关联的维修工单文本）}}
\codeblank{}
\codeline{\textcolor{cmtgray}{\#\ 变体：层级社区摘要（Microsoft\ GraphRAG,\ 2024）}}
\codeline{\textcolor{cmtgray}{\#\ \ \ 离线对图做社区检测，逐层生成社区摘要；}}
\codeline{\textcolor{cmtgray}{\#\ \ \ 全局性问题（"本月质量问题集中在哪类设备？"）检索社区摘要而非节点}}
\codeblank{}
\end{codebox}

\assetsource{ch08}{graphrag-pattern}

结构性问题（关联、路径、聚合）走图，描述性知识（维修案例的文字经过）走向量，
两路结果合并注入。检索的「锚定 \(\rightarrow\) 扩展 \(\rightarrow\) 挂载文档」
顺序保证了向量检索是在图谱圈定的范围内进行的，相关性显著好于全库盲检。
对「本月质量问题集中在哪类设备」这类全局性问题，
还可以引入层级社区摘要的变体：离线对图做社区检测、逐层生成摘要，
查询时检索摘要而非节点——代价是离线计算与摘要更新的运维成本，
适合图谱规模与查询量都到位之后再上。

评测上，GraphRAG 的优势集中在多跳问题：构造 50--100 个两到三跳的问题对照测试，
单跳问题两种方案通常接近，两跳以上结构化检索的完整率优势会拉开——
这也提示我们按问题分布决定是否值得引入图检索，而不是为了技术而技术。

\subsection{何时查询、何时检索：与 Text2SPARQL 的协同}

Text2SPARQL 与 GraphRAG 都从图谱取事实，边界在哪里？
判据是\textbf{问题的答案形态}。
答案是确定的记录或聚合值（状态、参数、数量、清单）——走查询：
精确、可校验、成本最低。
答案需要综合多个事实并组织成叙述（影响分析、原因解释、方案比较）——
走 GraphRAG：子图提供素材，模型负责综合。
答案还需要非结构化佐证（维修经过、客诉原文）——GraphRAG 混合检索。
落成决策树就是三个问句：能写成一条 SELECT/ASK 吗？
需要跨多个实体综合吗？需要文档原文吗？
实践中两者共享实体链接与 Schema 组件，是同一座地基上的两扇门，
不存在二选一的架构决策。

\subsection{实施清单与成本核算}

GraphRAG 上线前过一遍清单：
锚定组件是否复用了别名表；扩展是否启用了关系白名单与类型过滤；
k 值与子图节点上限是否设了硬顶；序列化是否带方向说明；
答案是否回流第三层做断言校验。
成本核算抓两头：\textbf{延迟}（图查询通常几十毫秒，
序列化与模型生成占大头，端到端预算按交互场景定在二到五秒）；
\textbf{token}（子图序列化是上下文的主要增量，
按「平均每节点二十到四十 token」估算，上限节点数直接决定单次成本——
这也是裁剪排序必须认真做的经济学理由）。

\section{本体引导的 Agent：从「答错话」到「不做错事」}

前几节的系统回答问题；智能体（Agent）则\textbf{执行动作}——
改设备状态、派工单、下指令。风险随之升级：
问答系统的幻觉产出一句错话，Agent 的幻觉直接改写业务系统乃至物理世界。
「把 3 号车床转入维护」如果被执行在一台正在加工高价值工件的设备上，
损失立刻物化。动作型系统的防线必须比问答更硬。

\subsection{护栏协议：提议—校验—执行—审计}

本体护栏的协议把一次动作拆成四个阶段。
\textbf{提议}（propose）：LLM 把自然语言指令翻译成结构化动作
（动作名、目标实体、参数）——注意它只有提议权，没有执行权。
\textbf{校验}（validate）：本体对提议做形式检查，三道关卡递进——
动作名必须在本体定义的动作集内（拦词汇幻觉）、
目标实体必须存在于知识库（拦实体幻觉）、
动作必须满足公理约束（拦状态幻觉，见下）。
\textbf{执行}（execute）：只有 APPROVED 的提议才到达执行器，
执行器本身不做判断，职责单一。
\textbf{审计}（audit）：每个提议无论批准与否都写入审计日志，
拒绝的记录连同拒绝理由与证据一起保存——它们是改进系统的原料。

\subsection{状态机公理：动作约束的通用模式}

三道关卡中最有领域含金量的是公理校验。
制造场景里最典型的公理是\textbf{状态机}：
设备状态在 \texttt{Running}、\texttt{Idle}、\texttt{Maintenance}、
\texttt{Fault} 之间的合法转换是有限集合——
故障设备必须先维护才能开工（\texttt{Fault} \(\rightarrow\) \texttt{Running} 非法），
运行中的设备不能直接派新任务。
把转换集合写进本体而不是写死在代码里，得到三重好处：
约束是\textbf{声明式}的（增删一条转换不用改代码）、
\textbf{可审计}的（合规人员能直接读懂规则全集）、
\textbf{可复用}的（同一套公理同时服务问答校验与动作校验）。
这个模式不限于设备：订单状态流转、审批链、配方变更窗口，
凡是「哪些转换合法」的知识都适用。

\subsection{参考代码与运行合同的差别}

ch08 包保留了一份约三十行的历史参考实现，用来阅读三道校验的控制流：

\begin{codebox}
\codeline{\hspace*{4\ttcw}def\ validate(self,\ proposal:\ dict)\ ->\ AuditRecord:}
\codeline{\hspace*{8\ttcw}"""三道检查：词汇\ ->\ 实体\ ->\ 公理（与幻觉控制三层一一对应）"""}
\codeline{\hspace*{8\ttcw}action\ =\ proposal.get("action",\ "")}
\codeline{\hspace*{8\ttcw}target\ =\ proposal.get("target",\ "")}
\codeblank{}
\codeline{\hspace*{8\ttcw}\textcolor{cmtgray}{\#\ 检查1\ 词汇约束：动作必须在本体定义的动作集中}}
\codeline{\hspace*{8\ttcw}if\ action\ not\ in\ VALID\_ACTIONS:}
\codeline{\hspace*{12\ttcw}return\ self.\_audit(action,\ target,\ "REJECTED",}
\codeline{\hspace*{31\ttcw}f"未知动作\ '\{action\}'（词汇幻觉）")}
\codeblank{}
\codeline{\hspace*{8\ttcw}\textcolor{cmtgray}{\#\ 检查2\ 实体存在性：目标必须是ABox中的个体}}
\codeline{\hspace*{8\ttcw}equip\ =\ self.onto["equipment"].get(target)}
\codeline{\hspace*{8\ttcw}if\ equip\ is\ None:}
\codeline{\hspace*{12\ttcw}return\ self.\_audit(action,\ target,\ "REJECTED",}
\codeline{\hspace*{31\ttcw}f"实体\ '\{target\}'\ 不存在于知识库（实体幻觉）")}
\codeblank{}
\codeline{\hspace*{8\ttcw}\textcolor{cmtgray}{\#\ 检查3\ 公理约束：状态机与前置条件}}
\codeline{\hspace*{8\ttcw}cur\ =\ equip["status"]}
\codeline{\hspace*{8\ttcw}if\ action\ ==\ "set\_status":}
\codeline{\hspace*{12\ttcw}new\ =\ proposal.get("value",\ "")}
\codeline{\hspace*{12\ttcw}if\ (cur,\ new)\ not\ in\ self.onto["status\_transitions"]:}
\codeline{\hspace*{16\ttcw}return\ self.\_audit(action,\ target,\ "REJECTED",}
\codeline{\hspace*{35\ttcw}f"非法状态转换\ \{cur\}\ ->\ \{new\}（违反状态机公理）",}
\codeline{\hspace*{35\ttcw}\{"current\_status":\ cur\})}
\codeline{\hspace*{8\ttcw}pre\ =\ self.onto["action\_preconditions"].get(action)}
\codeline{\hspace*{8\ttcw}if\ pre\ and\ cur\ not\ in\ pre["allowed\_status"]:}
\codeline{\hspace*{12\ttcw}return\ self.\_audit(action,\ target,\ "REJECTED",}
\codeline{\hspace*{31\ttcw}f"前置条件不满足：\{action\}\ 要求状态\ "}
\codeline{\hspace*{31\ttcw}f"\{pre['allowed\_status']\}，当前为\ \{cur\}",}
\codeline{\hspace*{31\ttcw}\{"current\_status":\ cur\})}
\codeblank{}
\codeline{\hspace*{8\ttcw}return\ self.\_audit(action,\ target,\ "APPROVED",\ "通过全部本体校验",}
\codeline{\hspace*{27\ttcw}\{"current\_status":\ cur\})}
\end{codebox}
\color{black}

\assetsource{ch08}{ontology-guided-agent}

逐段读这个 \texttt{validate} 方法，能看到三类幻觉与三道检查的一一对应；但该 asset
的 role 是 \texttt{legacy\_case\_reference}，不是原生 runner：
第一段拦「词汇幻觉」——LLM 发明的动作名（比如「超频运行」）在
\texttt{VALID\_ACTIONS} 里查无此项，直接 REJECTED；
第二段拦「实体幻觉」——目标设备在 ABox 中不存在，说明实体是编造的；
第三段拦「状态幻觉」——动作本身合法、实体也存在，
但当前状态不允许这个转换。历史代码用字典模拟状态机，只表达控制流；
当前 package 尚未把它替换为绑定 snapshot、查询、权限与
动作执行器的端到端合同。
每个分支都以 \texttt{\_audit} 收尾，确保审计日志的完整性不依赖调用方自觉。

scenario registry 记录了三个\textbf{参考 verdict}，而不是建议读者直接运行 Python：
「把 3 号车床转入维护」——当前状态 Running，转换合法，放行并执行；
「让 7 号铣床直接开工」——它处于 Fault 状态，
状态机公理拒绝 \texttt{Fault} \(\rightarrow\) \texttt{Running}，阻断；
「给幻影设备派活」——\texttt{Ghost\_999} 不存在于知识库，阻断。
这些预期结果只有在模型、prompt、ontology snapshot、权限、执行器与 receipt
全部绑定后，才可升级为原生执行证据。当前 runner 对未支持操作合同 fail closed。

\subsection{审计日志的字段设计与保留}

审计要经得起「半年后回查」的考验，最小字段集包括：
时间戳（UTC，精确到毫秒）、会话与请求标识、代理的用户身份与角色、
原始自然语言指令、结构化提议全文、裁决结果与命中的具体规则
（哪条公理、哪个前置条件）、证据快照（裁决时刻相关实体的状态）、
执行结果或阻断说明。其中\textbf{证据快照}最容易被省略也最不该省略——
知识库是活的，半年后那台设备的状态早变了，
没有快照就无法重建「当时为什么这么判」。
保留策略按风险分级：查询类九十天滚动即可，
动作类与被拒绝的提议建议按行业合规要求长期归档
（它们分别是事故回溯与攻击取证的一手材料）。
格式上，审计记录本身就适合存成 RDF（PROV-O 词汇），
让「审计查询」也能用 SPARQL 表达——审计员问
「上月所有被状态机公理拦截的提议」就是一条查询的事。

\subsection{多 Agent 协作下的护栏}

复杂任务往往由多个 Agent 分工完成：计划 Agent 拆解任务、
执行 Agent 各管一摊（查询、排产、通知）。多 Agent 架构下护栏有两条部署原则。
\textbf{校验跟着动作走}：护栏部署在每个「能触碰真实系统」的执行 Agent 前，
而不是只在计划 Agent 后——计划阶段的描述性输出可以宽松，
落地为动作的那一刻必须过闸。
\textbf{冲突仲裁用本体}：两个 Agent 同时给一台设备派任务时，
冲突的发现不靠 Agent 之间互相通知（它们可能根本不知道彼此），
而靠共享知识库上的公理——第二个写入的任务断言会触发
「同设备时间重叠」冲突（第 5 章的 SWRL 规则原样复用），
仲裁逻辑收敛在一个地方。
这体现共享语义权威层的架构价值，但“有公理”还不够：冲突检查必须是 Semantica
已支持并绑定的场景，两个动作还要进入同一事务/并发控制与授权面。当前 ch08 并未
完成该操作合同，因此这里是设计规格，不是已运行保证。

\subsection{与工具调用协议的关系，以及人的位置}

当前主流的 Agent 框架都提供工具调用协议（让模型以结构化方式调用外部函数）。
本体护栏与这些协议是\textbf{正交}的：协议解决「模型如何表达动作」，
护栏解决「动作是否应该被执行」。无论底层用哪种协议，
护栏都作为独立的一层插在「模型输出」与「真实执行器」之间，
这种与协议无关的定位让它可以平移到任何 Agent 框架上。

人在环路（human-in-the-loop）是护栏的自然延伸：
被 REJECTED 的动作不是简单丢弃，而是转人工确认队列——
操作员看到「Agent 想让故障铣床开工，已被状态机公理拦截」，
既能纠正 Agent 的误解，也可能发现知识库本身过期（设备其实已修好）。
随着信任积累，低风险动作可以进入自动放行白名单，
高风险动作（影响安全、超过金额阈值）永久保留人工确认——
自动化程度是一个随运行证据演化的旋钮，而不是一次性的决定。

\subsection{权限模型：谁的提议、动谁的设备}

状态机公理回答「这个转换是否合法」，权限回答「这个主体是否有权提议它」。
两者都适合放进本体：给 Agent（以及它代理的用户）建立角色个体，
动作与角色之间用「可提议」关系连接，校验时多查一跳即可。
一个实用的最小权限矩阵：查询类动作全员开放；
状态变更类按设备归属限定到对应车间的角色；
跨车间调度、参数修改这类高影响动作只授予计划员角色且强制人工复核。
把矩阵放在本体里的收益与状态机相同——
权限变更是数据变更而非代码发布，审计时一条 SPARQL 就能导出全量授权清单。

\subsection{回滚与补偿：护栏之后的最后一米}

即使提议通过了全部校验，执行仍可能出错（执行器超时、下游系统拒绝）或
事后被证明不当（人工复盘发现误操作）。动作系统因此需要\textbf{补偿设计}：
每个动作登记其逆动作（转入维护 \(\leftrightarrow\) 恢复待机；
派工 \(\leftrightarrow\) 撤工单），审计日志记录到足以重建动作前状态的粒度。
注意不是所有动作都可逆——已经开始切削的工件回不去毛坯——
不可逆动作应在本体中显式标记，校验层对它们自动提升确认等级。
「可逆性」本身成为一条领域知识，这是本体方法论的一贯风格：
把工程判断显式化、形式化，而不是散落在工程师的直觉里。

\section{反向赋能：用大模型加速本体工程}

前面都在讲本体管住大模型；方向反过来同样成立——
本体工程中最耗时的环节，恰好是大模型擅长的语言任务。

\textbf{起草类层次}：把设备台账的类型列、工艺文件的术语表喂给模型，
让它提议初版类层次与属性清单。模型给出的草稿通常覆盖面好、
组织合理，但会夹带过度泛化的类与似是而非的层级——
这正是第 3 章 OntoClean 体检要抓的问题。
\textbf{抽取候选公理}：让模型从工艺规程文本中提取「数控车床必须配备冷却系统」
这类约束语句并翻译成候选公理，专家逐条裁决。
\textbf{生成 CQ 草稿}：把岗位手册或客服记录喂给模型，
让它归纳「这个领域的人会问什么」，产出能力问题候选清单——
访谈前先有一份草稿，专家的时间用在勘误而不是从零开始。

\subsection{哪些工序不适合交给模型}

并非本体工程的每个环节都欢迎自动化，三类决策应当明确保留给人。
\textbf{顶层区分}：物理对象与过程、类型与角色这类顶层切分
（第 2、3 章反复强调的根基）一旦定错，返工成本是全局性的；
模型对这类哲学性区分的提案看似合理实则随机，不值得冒险。
\textbf{同一性判据}：「序列号还是资产编号作为设备的同一性判据」
取决于企业的数据治理现实（哪个系统是权威源、编号会不会重用），
这是组织知识而非语言知识，模型无从知晓。
\textbf{命名与归属的政治决策}：「这个概念归生产部还是质量部定义」
「用集团标准词还是事业部惯用词」——这些决策的本质是利益协调，
交给模型不是效率问题而是责任错位。
一条简单的划分线：\textbf{语言密集、判断标准客观的环节交给模型起草；
后果重大、判断标准在组织内部的环节由人决策}。

\subsection{提示词也是资产}

反向赋能用到的提案提示词（「请从以下规程中提取约束语句并翻译为候选公理……」）
不应散落在工程师的聊天记录里。它们与查询模板、话术模板一样进版本库：
与本体版本联动（本体的类层次变了，起草提示词里的示例也要跟着变）、
带评测样例（同一份输入文档，新旧提示词的提案质量可对比）、
有负责人与变更记录。一个成熟团队的「提示词库」最终会与「本体库」「模板库」
并列为三大知识资产——它们共同构成组织在 AI 时代的工程沉淀，
而不是某次项目的临时产物。

三种用法共享同一条纪律：\textbf{LLM 提案、专家裁决、版本库收口}。
模型产出一律标记为草稿（draft），不经领域专家确认不得进入正式本体；
裁决记录（谁、何时、采纳或驳回、理由）随版本库保存。
这个流程之所以是当前最佳实践，在于它把双方的优势拼在了一起：
模型提供覆盖面与速度，专家提供判断与责任，本体的质量门槛没有降低，
而到达门槛的速度提高了。
需要警惕的两类风险也来自模型的统计本性：
\textbf{隐性偏置}（训练语料中的行业惯例未必符合本厂实际）
与\textbf{过度泛化}（把「常见搭配」当成「必然约束」提案成公理）。
对策无他：专家裁决环节不可省略，且对模型提案保持比对人工提案更高的怀疑度。

\subsection{一个端到端的工作流示例}

把三种用法串成一个真实形态的工作流：车间要为「热处理工段」扩展本体。
第一步，工程师把约五十页的热处理工艺规程与设备说明书交给模型，
要求输出术语候选表（约两百条，含出现频次与例句）；
专家用一小时圈出八十条有效术语——这一步替代了过去两周的人工通读。
第二步，模型基于有效术语提议类层次草稿与三十条候选公理
（「渗碳炉属于热处理设备」「淬火工序必须紧随加热工序」），
每条附带其来源段落；专家逐条裁决，采纳二十一条、修改六条、驳回三条——
驳回的一条正是「过度泛化」的典型：模型把某型号炉子的限制提案成了全类约束。
第三步，模型从规程的「常见问题」章节归纳出十二个 CQ 草稿，
专家补充三个、合并两个，形成该工段的验收基线。
全程模型产出都带 draft 标记进版本库，裁决记录可追溯。
两周的扩展工作压缩到三天，而且——更重要的是——
质量门槛（专家裁决、OntoClean 体检、CQ 验收）一个都没有少。

\section{评测体系：证明系统真的可信}

防线建完了，最后一个问题是：怎么向自己和监管方证明它有效？

\subsection{三个核心指标}

\begin{center}
\begin{tabularx}{\linewidth}{@{}lXX@{}}
\toprule
\textbf{指标} & \textbf{计算口径} & \textbf{要点} \\
\midrule
事实准确率 & 回答中可校验断言的正确数 / 可校验断言总数 & 以图谱为判据逐条核验 \\
拒答正确率 & 知识库确无记录时正确回答「无记录」的比例 & 防止系统「硬答」 \\
幻觉漏过率 & 流出闸门的编造内容数 / 总回答数 & 核心安全指标，目标趋零 \\
\bottomrule
\end{tabularx}
\end{center}

三个指标必须一起看：只盯准确率会鼓励系统在不确定时硬答；
只盯拒答率会让系统动辄拒绝、失去可用性；
漏过率则是安全底线，它的分子需要人工抽样审计来发现——
被闸门拦下的幻觉有日志可数，漏出去的只能靠人查。

\subsection{评测集与持续评测}

评测集按三类样本配比构造：\textbf{正例}（知识库中有明确答案的问题）、
\textbf{无记录例}（实体存在但所询事实无记录——专测拒答行为）、
\textbf{边界例}（实体不存在、问法含混、多跳组合——专测系统在模糊地带的表现）。
三类的配比建议向后两类倾斜，因为它们才是幻觉的高发区。
上线不是评测的终点：每次本体或提示词变更后跑回归集；
生产流量定期抽样人工审计；监控「拒答率」「Schema 校验失败率」等过程指标的漂移——
漂移往往先于事故出现，是免费的预警信号。

\subsection{上线门槛：一张 go/no-go 清单}

把本章的要求收拢成上线评审的检查清单，逐项有证据才放行：
图谱端点只读且操作白名单生效（演示一次注入被拦截）；
高频意图的模板覆盖率达标且全部通过验收查询；
生成路线四道校验齐备（各演示一个被拦截的反例）；
出口校验四项可独立开关且默认全开；
拒答与无记录话术经业务方评审；
评测集三类样本配比合理、基线指标达到事先约定的门槛；
红队首轮测试完成且发现项已闭环；
审计日志可按会话回放完整决策链；
回滚预案演练过一次（含不可逆动作的升级确认流程）。
清单的价值不在条目本身，而在它把「系统可信」从一句口号
变成了一组可验证的工程事实——这正是本书从第 1 章贯穿至此的方法论。

\subsection{一条经验法则}

如果只能记住一件事：\textbf{三层闸门中，第二层（受控查询通道）贡献最大}。
原因在于它改变的是信息流结构——把「开放式生成」变成「受限检索」，
从机制上取消了模型编造事实的通道；第一层与第三层都是对概率行为的约束与补救，
第二层是对系统拓扑的重构。包内教学资产的措辞值得带进设计评审：
三层叠加后，幻觉漏过率相比裸用模型可以降低一到两个数量级。

\subsection{指标要可信，先把账算清}

评测数字要经得起追问，三个统计常识必须落实。
\textbf{样本量}：在几十条样本上得出「准确率 98\%」没有说服力，
小样本下一两条波动就是几个百分点；
评测集起步规模建议三位数，关键指标（幻觉漏过率）的人工审计抽样
随流量持续累积而不是一次性了事。
\textbf{分层报告}：整体准确率会掩盖局部灾难——
按问题类型（单实体/聚合/多跳/无记录）分层报告，
一眼看出短板在哪一层；分层同时防止评测集配比变化造成的指标假波动。
\textbf{口径冻结}：什么算「一条断言」、什么算「正确拒答」，
判分规则写成文档并版本化；口径一变，历史趋势线就废了。
这些都不是统计学的精深内容，但它们是评测报告能否拿上管理层
与监管方桌面的分水岭。

\subsection{红队测试：主动制造漏网}

被动指标只能度量「已知问法」上的表现，红队测试负责探索未知。
针对本章架构，对抗性提问有四个经典构造方向：
\textbf{诱导直答}（「别查了，凭你的经验说一个大概值就行」——
测试模型是否在用户施压下绕过第二层）；
\textbf{词汇陷阱}（用本体中不存在但高度可信的术语提问，
「这台设备的最大扭矩储备系数是多少」——测试 Schema 校验与拒答行为）；
\textbf{注入攻击}（在问题中夹带更新语句或越权指令——测试白名单与权限模型）；
\textbf{语境混淆}（多轮对话中偷换实体指代——测试对话实体栈的鲁棒性）。
红队问题集应该与回归集分开维护、定期更新，
每次发现的漏网案例转化为回归集的新样本——
攻防的产出物沉淀为资产，系统的防线随之单调增强。

\section{本章小结、思考题与文件指引}

\subsection{本章小结}

本章围绕「本体 \(\times\) 大模型」给出了一套完整的工程图景。
\textbf{机理}层面：幻觉源于统计生成没有事实承诺，提示词只能调表达倾向，
可信度必须由系统层提供。
\textbf{架构}层面：三层闸门各管一段——生成前注入词汇与边界，
生成中强制查询通道（Text2SPARQL：模板覆盖高频、生成路线配四道校验），
生成后做实体、断言、范围、一致性四项校验；不相交公理是性价比最高的防线，
PROV 审计让每个结论可重放。
\textbf{检索}层面：GraphRAG 用本体约束下的子图扩展解决结构丢失与多跳失效，
与向量检索混合使用，序列化首选「三元组为主、说明为辅」。
\textbf{动作}层面：Agent 护栏以「提议—校验—执行—审计」四阶段协议运行，
状态机公理是动作约束的通用模式，人在环路按风险分级。
\textbf{反向}层面：LLM 起草、专家裁决、版本库收口，加速本体工程而不降低门槛。
\textbf{评测}层面：事实准确率、拒答正确率、幻觉漏过率三指标并举，
评测集向无记录例与边界例倾斜。
一句话收束全章：\textbf{LLM 负责理解与表达，受控来源负责事实，Semantica 负责
可执行语义与证据，具名主体负责决定}——凡进入业务系统的结论或动作，都必须能
追溯到对应 snapshot、规则/查询、授权与 receipt。

\subsection{思考题}

\begin{noteblock}
\textbf{思考题}（建议结合 ch08 package assets 与 manifest 作答）

1. 「张冠李戴」型幻觉为什么无法靠扩充训练语料根治？三层闸门中哪一层对它最有效？
（提示：从拼接发生的环节想。）

2. 设计第一层闸门时，如果本体太大无法全部注入上下文，你会按什么策略选取子集？
（提示：与用户问题的实体链接结果有关。）

3. 为什么说第二层闸门「改变了信息流的拓扑」？画出有无第二层时的两条信息流路径。
（提示：注意事实的来源节点变化。）

4. 路线B的四道校验中，去掉 Schema 校验会发生什么？构造一个会漏过的具体查询例子。
（提示：发明一个本体中不存在但听起来合理的属性名。）

5. 「P001 延期影响哪些客户订单」需要三跳。写出 GraphRAG 子图扩展时应使用的
关系白名单，并说明每个关系来自本体的哪条属性定义。
（提示：producedBy / partOf / orderedBy。）

6. Agent 护栏的三道检查分别拦截哪类幻觉？如果把检查顺序颠倒（先查状态机再查实体存在），
会出什么问题？（提示：对不存在的实体查状态会怎样。）

7. 状态机公理写在本体里与写死在代码里相比有哪三重好处？
各举一个会体现差异的运维场景。（提示：变更、审计、复用。）

8. 你的系统幻觉漏过率为 0，这是好消息吗？还需要检查哪两个指标才能下结论？
（提示：拒答率过高同样不可用。）
\end{noteblock}

\subsection{本章包资产与状态指引}

下列 asset 位于 \texttt{semantica.chapter\_packages.vol1.ch08}。

\begin{center}
\begin{tabularx}{\linewidth}{@{}>{\ttfamily\footnotesize}p{0.38\linewidth}X@{}}
\toprule
\normalfont\textbf{Semantica asset} & \textbf{内容与阅读建议} \\
\midrule
hallucination-control & 三层架构、类型拦截与评测指标；工程规则材料 \\
text2sparql & 三条路线、四道校验与失败回退；当前没有受 schema/CQ 限定的原生场景 \\
graphrag-pattern & 子图扩展、序列化与混合检索；教学 case \\
ontology-guided-agent & legacy case reference；操作合同 unsupported，不是可运行入口 \\
OE-V1-CH08-SCN-GUARDRAIL-001 & 参考 verdict 已登记；runner status 仍 blocked \\
\bottomrule
\end{tabularx}
\end{center}

读完本章，全书的概念旅程只剩最后一步：把第 1--8 章已经支持的能力与仍然阻断的
能力一起纳入一个诚实的 package/release 视图。这就是第 9 章的任务。
